
\chapter{סיכום: לאן מוביל עתיד התיאום}

\section{תמונת המצב}

תחום תיאום כלים מרובים עבר דרך ארוכה תוך שלוש שנים: מ-\textenglish{ReAct} הפשוט \cite{yao2022react} לצינורות היברידיים מרובי-סוכנים עם ניהול עלויות, אבטחה ודיבייט אוטומטי. הממצאים המרכזיים:

\begin{itemize}
\item ניתוב חכם (\textenglish{Router-Skill}) חוסך עד \textenglish{23\%} בעלות אסימון.
\item מקביליות מפחיתה \textenglish{latency} ב-\textenglish{$\sim$6×} עבור \textenglish{16} סוכנים.
\item \textenglish{Watchdog} ומגבלות איטרציה קריטיים לייצוב צינורות ייצור.
\end{itemize}

\section{כיוונים עתידיים}

העתיד מחייב: (א) פרוטוקולי \textenglish{A2A} (\textenglish{Agent-to-Agent}) תקניים לאינטרופרביליות; (ב) ערבויות פורמליות על עלות אסימון (\textenglish{token budget guarantees}); (ג) שכבות אבטחה מובנות נגד \textenglish{prompt injection} בכלל המסגרות \cite{nakash2024skillsieve}. האינטגרציה עם \textenglish{SkillSieve} \cite{nakash2024skillsieve} מהווה צעד מוצלח לכיוון מסגרות סוכנים עמידות לאיומי שרשרת אספקה.
